\section{The Picard group}\label{sec:pic}
% ===============================================================

Divisors do not only govern functions; they also produce line bundles, and the
group of line bundles is the geometric counterpart of the class group. In this
section we compute it. We work throughout in the holomorphic (analytic) category:
$\PP$ is a Riemann surface and all bundles, sections and transition functions are
holomorphic. See Remark~\ref{rem:algebraic} for the algebraic counterpart.

\subsection{Line bundles by transition functions}

\begin{definition}
A \emph{holomorphic line bundle} on a Riemann surface $X$ is a holomorphic vector
bundle $L \to X$ of rank $1$. The set $\Pic(X)$ of isomorphism classes of
holomorphic line bundles is a group under tensor product: the identity is the
class of the trivial bundle $\Ostr = X \times \CC$, and the inverse of $[L]$ is
the class of the dual bundle $L^{\vee}$, since $L \otimes L^{\vee} \cong \Ostr$.
\end{definition}

Concretely, a line bundle is described by trivialising it over an open cover and
recording how the trivialisations differ. We use only the standard cover
$\PP = U_0 \cup U_1$ of Section~\ref{sec:sphere}, on whose overlap
$U_0 \cap U_1 = \CC^\times$ we have the coordinate $z$.

\begin{definition}\label{def:cocycle}
Let $\Ostr^\ast(V)$ denote the group, under multiplication, of nowhere-vanishing
holomorphic functions on an open $V \subseteq \PP$. Given
$g \in \Ostr^\ast(\CC^\times)$, let $L_g$ be the line bundle obtained by glueing
$U_0 \times \CC$ to $U_1 \times \CC$ along the overlap via
$(z, \xi_1) \mapsto (z, g(z)\xi_1)$. We call $g$ the \emph{transition function}
of $L_g$ relative to this cover.
\end{definition}

\begin{smjlemma}\label{lem:cocycleiso}
$L_g \cong L_{g'}$ if and only if there are $h_0 \in \Ostr^\ast(U_0)$ and
$h_1 \in \Ostr^\ast(U_1)$ with $g' = h_0\, g\, h_1^{-1}$ on $\CC^\times$.
Moreover $L_g \otimes L_{g'} \cong L_{gg'}$.
\end{smjlemma}

\begin{proof}
An isomorphism $L_g \to L_{g'}$ is precisely a pair of nowhere-vanishing
holomorphic functions $h_i$ on $U_i$ ($i = 0,1$) --- the two trivialised
descriptions of the isomorphism --- compatible on the overlap, and compatibility
reads $g' = h_0 g h_1^{-1}$. The tensor statement holds because tensoring
multiplies transition functions.
\end{proof}

We shall use one imported fact, the only result in this paper not proved here.

\begin{fact}[{\cite[\S30]{Forster}}]\label{fact:trivialise}
Every holomorphic line bundle on $\PP$ admits trivialisations over $U_0$ and over
$U_1$, and is therefore isomorphic to $L_g$ for some
$g \in \Ostr^\ast(\CC^\times)$.
\end{fact}

This holds because each $U_i$ is biholomorphic to $\CC$, and every holomorphic
line bundle on $\CC$ is trivial. We isolate it as an explicit hypothesis rather
than absorbing it silently.

\subsection{Winding numbers and the classification of transition functions}

\begin{definition}
For $g \in \Ostr^\ast(\CC^\times)$ the \emph{winding number} is
\[
\wind(g) := \frac{1}{2\pi i}\oint_{|z|=1} \frac{g'(z)}{g(z)}\,dz .
\]
\end{definition}

It is an integer --- the winding number of the loop $g(e^{i\theta})$ about the
origin --- and $\wind(gh) = \wind(g) + \wind(h)$, since $(gh)'/(gh) = g'/g +
h'/h$. Note $\wind(z^n) = n$.

\begin{smjlemma}\label{lem:windcoboundary}
If $h_0 \in \Ostr^\ast(U_0)$ and $h_1 \in \Ostr^\ast(U_1)$ then
$\wind(h_0) = \wind(h_1) = 0$.
\end{smjlemma}

\begin{proof}
$h_0$ is holomorphic and nowhere zero on $U_0 = \CC$, so $h_0'/h_0$ is
holomorphic on the disc $|z| \leq 1$ and its integral over the boundary vanishes
by Cauchy's theorem. For $h_1$, pass to the coordinate $w = 1/z$: the function
$w \mapsto h_1(1/w)$ is holomorphic and nowhere zero on $U_1 \cong \CC$, so by the
same argument $\oint_{|w|=1}\frac{d}{dw}\log h_1 \,dw = 0$; substituting
$w = 1/z$ changes the integral by a sign, so $\wind(h_1) = 0$ as well.
\end{proof}

\begin{smjlemma}[Factorisation of transition functions]\label{lem:cousin}
Every $g \in \Ostr^\ast(\CC^\times)$ can be written as
\[
g(z) = z^{\,n}\, h_0(z)\, h_1(z)^{-1}, \qquad n = \wind(g),
\]
with $h_0 \in \Ostr^\ast(U_0)$ and $h_1 \in \Ostr^\ast(U_1)$.
\end{smjlemma}

\begin{proof}
Put $n := \wind(g)$ and $G(z) := z^{-n}g(z)$, so that $\wind(G) = \wind(g) - n =
0$. The function $G'/G$ is holomorphic on $\CC^\times$, and $\CC^\times$ has the
homotopy type of a circle, so $G'/G$ admits a primitive on $\CC^\times$ precisely
when its integral over $|z| = 1$ vanishes \cite[Ch.~4]{Ahlfors} --- which is
exactly the statement $\wind(G) = 0$. Let $\varphi$ be such a primitive. Then
$\big(Ge^{-\varphi}\big)' = e^{-\varphi}(G' - G\varphi') = 0$, so
$G = \alpha e^{\varphi}$ for a constant $\alpha \in \CC^\times$; absorbing
$\alpha$ into $\varphi$ (replace $\varphi$ by $\varphi + \log\alpha$ for any
choice of logarithm) we may assume $G = e^{\varphi}$.

Now expand $\varphi$ in a Laurent series on $\CC^\times$
(Theorem~\ref{thm:laurent}), $\varphi(z) = \sum_{k \in \ZZ} a_k z^k$, and split
it into its nonnegative and negative parts,
\[
\varphi_0(z) := \sum_{k \geq 0} a_k z^{k}, \qquad
\varphi_1(z) := \sum_{k < 0} a_k z^{k} = \sum_{j > 0} a_{-j} w^{j} .
\]
The series $\varphi_0$ converges on all of $\CC$, hence is holomorphic on $U_0$;
the series $\varphi_1$, viewed in the coordinate $w = 1/z$, is a power series in
$w$ with no constant term, converging on all of $\CC$, hence holomorphic on
$U_1$. On the overlap $\varphi = \varphi_0 + \varphi_1$. Setting
$h_0 := e^{\varphi_0} \in \Ostr^\ast(U_0)$ and
$h_1 := e^{-\varphi_1} \in \Ostr^\ast(U_1)$ gives
$G = e^{\varphi_0}e^{\varphi_1} = h_0 h_1^{-1}$ and therefore
$g = z^n h_0 h_1^{-1}$.
\end{proof}

\begin{theorem}\label{thm:picZ}
The map $\ZZ \to \Pic(\PP)$, $n \mapsto [L_{z^n}]$, is a group isomorphism.
\end{theorem}

\begin{proof}
It is a homomorphism by the tensor statement of Lemma~\ref{lem:cocycleiso}, since
$z^{m}z^{n} = z^{m+n}$.

\emph{Surjectivity.} By Fact~\ref{fact:trivialise} every line bundle is $L_g$ for
some $g \in \Ostr^\ast(\CC^\times)$. By Lemma~\ref{lem:cousin},
$g = z^{n}h_0h_1^{-1}$ with $n = \wind(g)$, so $L_g \cong L_{z^n}$ by
Lemma~\ref{lem:cocycleiso}.

\emph{Injectivity.} Suppose $L_{z^m} \cong L_{z^n}$. By
Lemma~\ref{lem:cocycleiso} there are $h_i \in \Ostr^\ast(U_i)$ with
$z^{m} = h_0 z^{n} h_1^{-1}$. Taking winding numbers and using
Lemma~\ref{lem:windcoboundary},
$m = \wind(h_0) + n - \wind(h_1) = n$.
\end{proof}

\subsection{The line bundle of a divisor}

\begin{definition}\label{def:OD}
Let $D \in \Div(\PP)$. For open $U \subseteq \PP$ set
\[
\Ostr(D)(U) := \{\, f \in \CC(z)^\times : \big(\divv(f) + D\big)\big|_U \geq 0 \,\}
\cup \{0\},
\]
using the notation of Definition~\ref{def:effective}. Explicitly, $f$ belongs to
$\Ostr(D)(U)$ if and only if $\ord_p(f) + n_p \geq 0$ for every $p \in U$, that
is, if $f$ is allowed a pole of order at most $n_p$ at $p$ when $n_p > 0$ and is
required to vanish to order at least $-n_p$ when $n_p < 0$. The zero function is
adjoined separately because $\divv(0)$ is undefined.
\end{definition}

\begin{smjproposition}\label{prop:ODbundle}
$\Ostr(D)$ is a line bundle: every $p \in \PP$ has a neighbourhood $V$ on which
$\Ostr(D)$ is free of rank one over the holomorphic functions, with generator
$t_p^{-n_p}$.
\end{smjproposition}

\begin{proof}
Choose $V \ni p$ meeting $\supp(D)$ only in $p$, possible since $\supp(D)$ is
finite. For $f$ meromorphic on $V$ the condition
$(\divv(f)+D)|_V \geq 0$ reads $\ord_p(f) \geq -n_p$ together with
$\ord_q(f) \geq 0$ for $q \in V\setminus\{p\}$; by
Propositions~\ref{prop:ord} and \ref{prop:ordhom} this holds if and only if
$t_p^{\,n_p}f$ is holomorphic on $V$. Hence
$\Ostr(D)(V) = t_p^{-n_p}\cdot \Ostr(V)$, as claimed.
\end{proof}

\begin{smjproposition}\label{prop:ODinvariant}
If $D \sim E$ then $\Ostr(D) \cong \Ostr(E)$. In particular, if $D$ is principal
then $\Ostr(D) \cong \Ostr$.
\end{smjproposition}

\begin{proof}
Let $D - E = \divv(f)$. For any open $U$ and any $g \in \Ostr(D)(U)$,
\[
\divv(gf) + E = \divv(g) + \divv(f) + E = \divv(g) + (D-E) + E = \divv(g) + D,
\]
which is $\geq 0$ on $U$ by hypothesis; hence multiplication by $f$ maps
$\Ostr(D)(U)$ into $\Ostr(E)(U)$. Multiplication by $f^{-1}$ is an inverse map,
so this is an isomorphism, compatible with restriction and with multiplication by
holomorphic functions. Taking $E = 0$ gives the last statement, since
$\Ostr(0) = \Ostr$.
\end{proof}

\begin{smjproposition}\label{prop:Oninfty}
For $n \in \ZZ$ we have $\Ostr(n(\infty)) \cong L_{z^n}$.
\end{smjproposition}

\begin{proof}
Apply Proposition~\ref{prop:ODbundle} with the cover $\{U_0, U_1\}$, which is
legitimate because $\supp(n(\infty)) = \{\infty\}$ meets $U_0$ not at all and
$U_1$ only in $\infty$. Over $U_0$ the condition is that $f$ be holomorphic, so
$\Ostr(n(\infty))(U_0) = \Ostr(U_0)$ with generator $e_0 = 1$. Over $U_1$ the
generator is $t_\infty^{-n} = z^{n}$, so $e_1 = z^n$.

A section $f$ over the overlap has coordinates $\xi_0, \xi_1$ determined by
$f = \xi_0 e_0 = \xi_1 e_1$, i.e.\ $\xi_0 = f$ and $\xi_1 = f z^{-n}$. Hence
$\xi_0 = z^{n}\xi_1$, so the transition function relative to $\{U_0, U_1\}$ is
$z^{n}$, and $\Ostr(n(\infty)) \cong L_{z^n}$ by Definition~\ref{def:cocycle}.
\end{proof}

\subsection{The Picard group of $\PP$}

\begin{theorem}\label{thm:pic}
The assignment $D \mapsto \Ostr(D)$ induces a group isomorphism
\[
\Cl(\PP) \xrightarrow{\ \sim\ } \Pic(\PP),
\]
and consequently $\Pic(\PP) \cong \ZZ$. Writing $\Ostr(n) := \Ostr(n(\infty))$,
every holomorphic line bundle on $\PP$ is isomorphic to $\Ostr(n)$ for a unique
$n \in \ZZ$, namely $n = \wind$ of its transition function.
\end{theorem}

\begin{proof}
The assignment is well defined on classes by Proposition~\ref{prop:ODinvariant}
and lands in $\Pic(\PP)$ by Proposition~\ref{prop:ODbundle}. Consider the
diagram of maps
\[
\Cl(\PP) \xrightarrow{\ [D]\mapsto[\Ostr(D)]\ } \Pic(\PP)
\xleftarrow{\ n \mapsto [L_{z^n}]\ } \ZZ .
\]
By Proposition~\ref{prop:normalform} every class satisfies $[D] =
[(\deg D)(\infty)]$, so by Proposition~\ref{prop:Oninfty} the left-hand map sends
$[D]$ to $[L_{z^{\deg D}}]$; that is, it equals the composite of
$\overline{\deg} : \Cl(\PP) \to \ZZ$ with $n \mapsto [L_{z^n}]$. The first factor
is an isomorphism by Theorem~\ref{thm:cl} and the second by
Theorem~\ref{thm:picZ}, so the left-hand map is an isomorphism as well. In
particular it is a homomorphism, and $\Pic(\PP) \cong \Cl(\PP) \cong \ZZ$. The
final statement is Theorem~\ref{thm:picZ} combined with
Proposition~\ref{prop:Oninfty}.
\end{proof}

Thus line bundles on $\PP$, like divisor classes, are classified by a single
integer, and that integer is the degree.

\begin{remark}[The algebraic counterpart]\label{rem:algebraic}
We have worked analytically throughout: $\Pic(\PP)$ was the group of
\emph{holomorphic} line bundles on a complex manifold. In algebraic geometry one
defines $\Pic$ as the group of invertible sheaves on a scheme, and proves
$\Cl(X) \cong \Pic(X)$ for a noetherian, integral, separated, locally factorial
scheme \cite[Cor.~II.6.16]{Hartshorne}; for $X = \mathbb{P}^1_{\CC}$ this gives
$\Pic(\mathbb{P}^1_\CC) \cong \ZZ$ as well \cite[Prop.~II.6.4]{Hartshorne},
\cite[Tag 0BXJ]{Stacks}. That the algebraic and analytic answers agree is not a
formality: it is an instance of Serre's GAGA theorems \cite{GAGA}, which compare
coherent sheaves on a projective variety with their analytic counterparts. We
have avoided relying on this comparison by computing the analytic group directly.
\end{remark}

% ===============================================================
